\section*{ДОДАТКИ}

\subsection*{Додаток A. Псевдокод алгоритмів (розділ 4)}

\begin{lstlisting}[caption={Нормалізація DOI (\texttt{normalize\_doi})},label={lst:a-normalize-doi}]
normalize_doi(doi):
    doi = lowercase(doi)
    doi = remove_prefixes(doi, ["https://doi.org/", "dx.doi.org/"])
    doi = rstrip(doi)
    assert matches_regex(doi, r"^10\.\d{4,9}/.+$")
\end{lstlisting}

\begin{lstlisting}[caption={Первинний матчинг за DOI (\texttt{primary\_match})},label={lst:a-primary-match}]
primary_match(new_pub, existing_pubs):
    IF new_pub.doi EXISTS AND is_valid(new_pub.doi):
        normalized = normalize_doi(new_pub.doi)
        FOR each existing IN existing_pubs:
            IF normalize_doi(existing.doi) == normalized:
                RETURN existing.index
    RETURN -1  // Немає збігу
\end{lstlisting}

\begin{lstlisting}[caption={Вторинний матчинг за Title+Year (\texttt{secondary\_match})},label={lst:a-secondary-match}]
secondary_match(new_pub, existing_pubs):
    normalized_title = normalize_title(new_pub.title)
    FOR each existing IN existing_pubs:
        IF similarity(normalized_title, normalize_title(existing.title)) > 0.9:
            IF new_pub.year == existing.year:
                RETURN existing.index
    RETURN -1
\end{lstlisting}

\begin{lstlisting}[caption={Нормалізація назви (\texttt{normalize\_title})},label={lst:a-normalize-title}]
normalize_title(title):
    title = lowercase(title)
    title = remove_punctuation(title)
    title = remove_stopwords(title)
    title = stem_or_lemmatize(title)
\end{lstlisting}

\begin{lstlisting}[caption={Злиття метаданих (\texttt{merge})},label={lst:a-merge}]
merge(existing, new_pub):
    // Citation count: максимум
    existing.citations = MAX(existing.citations, new_pub.citations)

    // DOI: заповнити якщо відсутній
    IF existing.doi IS NULL AND new_pub.doi EXISTS:
        existing.doi = new_pub.doi

    // Open Access: оптимістичний підхід
    IF new_pub.is_open_access == TRUE:
        existing.is_open_access = TRUE

    // Keywords: об'єднання унікальних
    existing.keywords = UNION(existing.keywords, new_pub.keywords)

    // Sources: трекінг походження
    existing.sources.ADD(new_pub.source)
\end{lstlisting}

\begin{lstlisting}[caption={Трекінг провенансу (структури \texttt{sources}/\texttt{last\_updated})},label={lst:a-provenance}]
sources: ['orcid', 'openalex', 'scopus']
last_updated: {
    'orcid': '2024-01-15',
    'openalex': '2024-01-20',
    'scopus': '2024-01-10'
}
\end{lstlisting}

\subsection*{Додаток B. Приклади коду та журнали виконання (розділ 5)}

\begin{lstlisting}[language=Java,caption={Декларація tools для MCP Function Calling (TypeScript)},label={lst:b-tools}]
const tools = [
  {
    name: "get_author_metrics",
    description: "Отримати детальні метрики для дослідника: h-index, \n                  кількість цитувань, тренди публікацій, найкращі роботи",
    parameters: {
      type: "object",
      properties: {
        orcid: {
          type: "string",
          description: "ORCID iD (наприклад, 0000-0003-0242-2234)"
        }
      },
      required: ["orcid"]
    }
  },
  {
    name: "search_scientific_works",
    description: "Пошук наукових публікацій за темою, \n                  ім’ям автора, установою або ключовими словами",
    parameters: {
      type: "object",
      properties: {
        query: { type: "string", description: "Пошуковий запит" },
        year_from: { type: "integer", description: "Початковий рік (фільтр)" },
        year_to: { type: "integer", description: "Кінцевий рік (фільтр)" },
        open_access: { type: "boolean", description: "Фільтр: лише OA" }
      },
      required: ["query"]
    }
  },
  {
    name: "compare_researchers",
    description: "Порівняти метрики кількох дослідників",
    parameters: {
      type: "object",
      properties: {
        orcids: {
          type: "array",
          items: { type: "string" },
          description: "Список ORCID iD для порівняння"
        }
      },
      required: ["orcids"]
    }
  },
  {
    name: "get_department_summary",
    description: "Агрегувати статистики для кафедри/підрозділу",
    parameters: {
      type: "object",
      properties: {
        researchers: {
          type: "array",
          items: { type: "string" },
          description: "Список ORCID iD членів підрозділу"
        },
        period: { type: "string", description: "Період (наприклад, '2020-2024')" }
      },
      required: ["researchers"]
    }
  },
  {
    name: "generate_report",
    description: "Згенерувати форматований звіт у вказаному форматі",
    parameters: {
      type: "object",
      properties: {
        type: {
          type: "string",
          enum: ["annual", "attestation", "grant", "custom"],
          description: "Тип звіту"
        },
        format: {
          type: "string",
          enum: ["docx", "xlsx", "pdf"],
          description: "Формат виводу"
        },
        template: { type: "string", description: "Назва шаблону (для custom)" }
      },
      required: ["type", "format"]
    }
  }
];
\end{lstlisting}

\begin{lstlisting}[caption={Лог виконання: крок 1 (збір з ORCID)},label={lst:b-orcid-log}]
Отримано 45 публікацій для дослідника X
Джерело: orcid
\end{lstlisting}

\begin{lstlisting}[caption={Лог виконання: крок 2 (збір з OpenAlex)},label={lst:b-openalex-log}]
Отримано 52 публікації для дослідника X
Джерело: openalex
Виявлено: 38 потенційних дублікатів за DOI
\end{lstlisting}

\begin{lstlisting}[caption={Лог виконання: крок 3 (злиття)},label={lst:b-merge-log}]
Зіставлено за DOI: 38 записів -> злиті
Зіставлено за Title+Year: 4 записи -> злиті
Нові унікальні з OpenAlex: 10 записів -> додано
Нові унікальні з ORCID: 3 записи -> додано
\end{lstlisting}

\begin{lstlisting}[caption={Лог виконання: результат формування Golden Record},label={lst:b-result-log}]
Golden Records: 55 унікальних публікацій
Відстежено джерела: {orcid: 45, openalex: 52}
Частка дублікатів до злиття: 23%
Частка дублікатів після злиття: 0%
\end{lstlisting}

\subsection*{Додаток В. Розрахунок статистики ефективності (відтворювана модель)}
\label{app:efficiency}

\textbf{Статус доказів.} Наведений нижче розрахунок отримано на \emph{синтетичній узагальненій моделі} факультету (без реальної установи й персональних даних): канонічні праці проєктуються у записи ORCID/OpenAlex/Scopus із контрольованою, gold-розміченою дуплікацією та спотворенням метаданих, після чого запускається \emph{реальний} двигун Golden Record (\texttt{services/goldenRecord}) і базові методи DOI-only та Title+Year на ідентичних входах. Оцінки ручної праці --- це \textbf{прозора аналітична модель із задекларованими коефіцієнтами, а не вимірювання за участю людей}. Тому показники слід читати як відтворювані \emph{модельні} оцінки, що потребують подальшої калібровки на реальних даних (\textit{Evidence gap}).

\textbf{Модель факультету.} \textbf{50} НПП у \textbf{4} кафедрах одного факультету; період 2020--2024; частка дублікатів $0{,}2$; множник спотворення $1{,}0$; дисциплінарний мікс STEM/соціогуманітарні/медико-біологічні $= 0{,}40/0{,}35/0{,}25$. Статистика --- метод Монте-Карло, $K=30$ незалежних сідів (детермінований ГВЧ \texttt{mulberry32}); усі значення подано як середнє $\pm$ 95\% ДІ.

\textbf{Модель ручної праці} (без системи), хв:
\begin{equation}
T_{\text{ручна}} = a\cdot N_{\text{записів}} + b\cdot N_{\text{дублікат-рішень}} + c\cdot N_{\text{звітів}},\quad a=2{,}5,\ b=5,\ c=45,
\end{equation}
де $a$ --- хв на пошук і верифікацію одного запису, $b$ --- хв на одне рішення про злиття дубліката, $c$ --- фіксований час підготовки/форматування звіту кафедри.

\begin{longtable}{@{}p{0.62\textwidth}p{0.32\textwidth}@{}}
\caption{Статистика ефективності для модельного факультету (50 НПП, 4 кафедри; середнє $\pm$ 95\% ДІ, $K=30$)}\label{tab:efficiency-50}\\
\toprule
\textbf{Показник} & \textbf{Значення} \\
\midrule
\endfirsthead
\toprule
\textbf{Показник} & \textbf{Значення} \\
\midrule
\endhead
Мультиджерельних записів до злиття, $N_{\text{записів}}$ & $\approx 1167$ \\
Дублікат-рішень (gold), $N_{\text{дублікат-рішень}}$ & $\approx 238$ \\
\midrule
Ручна праця --- збір і верифікація ($a\cdot N_{\text{записів}}$) & $48{,}6$ год \\
Ручна праця --- дедуплікація ($b\cdot N_{\text{дублікат-рішень}}$) & $19{,}9$ год \\
Ручна праця --- фіксоване форматування ($c\cdot N_{\text{звітів}}$) & $3{,}0$ год \\
\textbf{Разом ручна праця (модель)} & \textbf{$71{,}5$ год} ($4288\pm204$ хв) \\
\quad у т.\,ч. на один звіт кафедри & $17{,}9$ год \\
\midrule
\textbf{Автоматизовна частка} (збір + дедуплікація) & \textbf{$95{,}8\%$} \\
Час двигуна (дедуплікація + злиття, лише compute) & $0{,}98\pm0{,}21$ с \\
Залишкова HITL-перевірка: позначено пар на ревізію & $56{,}7$ пар/НПП \\
\bottomrule
\end{longtable}

\textbf{Якість, що робить автоматизацію довірливою.} Скорочення ручної праці виправдане лише за достатньої якості дедуплікації та злиття (табл.~\ref{tab:quality-50}). Golden Record дає суттєво вищий recall за DOI-only (бо ловить дублікати без DOI) за збереження високої precision.

\begin{longtable}{@{}p{0.26\textwidth}p{0.17\textwidth}p{0.17\textwidth}p{0.17\textwidth}p{0.15\textwidth}@{}}
\caption{Якість дедуплікації за методами (\%, середнє $\pm$ 95\% ДІ, $K=30$)}\label{tab:quality-50}\\
\toprule
\textbf{Метод} & \textbf{Precision} & \textbf{Recall} & \textbf{F1} & \textbf{B\textsuperscript{3} F1} \\
\midrule
\endfirsthead
\toprule
\textbf{Метод} & \textbf{Precision} & \textbf{Recall} & \textbf{F1} & \textbf{B\textsuperscript{3} F1} \\
\midrule
\endhead
DOI-only & $100\pm0$ & $42{,}4\pm1{,}4$ & $59{,}5\pm1{,}4$ & $93{,}7\pm0{,}2$ \\
Title+Year & $82{,}5\pm1{,}4$ & $38{,}9\pm1{,}3$ & $52{,}8\pm1{,}3$ & $92{,}6\pm0{,}2$ \\
\textbf{Golden Record} & $87{,}8\pm1{,}0$ & $67{,}2\pm1{,}3$ & $76{,}1\pm1{,}0$ & $95{,}7\pm0{,}2$ \\
\bottomrule
\end{longtable}

Польова повнота/точність злитого запису (Golden Record): DOI наявний $45{,}6\%$, DOI коректний $67{,}9\%$, рік $80{,}8\%$, назва $99{,}97\%$. (Окремо: $100\%$ україномовних назв були б «вихолощені» старим нормалізатором із багом для кирилиці --- демонстрація виправленої вади.)

\textbf{Інтерпретація.} Модель показує, що $\approx 95{,}8\%$ модельованої ручної праці (збір/верифікація + дедуплікація) система переводить у $\approx 1$~с детермінованого обчислення; залишковий людський внесок зводиться до HITL-перевірки позначених невпевнених пар і фінального контролю, тоді як фіксоване форматування звітів автоматизується шаблонами. Це не «нульова праця людини», а \emph{зсув} зусиль від $O(N_{\text{записів}})$ ручної консолідації до $O(\text{позначених})$ верифікації.
\iftrace
У термінах TRACE час двигуна відповідає метриці TCC (L3), HITL-навантаження --- RBI (L4), а автоматизовна частка є прямим свідченням принципу Model Parsimony (детерміноване ядро замість LLM для structured-задачі), що операціоналізується через CPR.
\fi

\textbf{Відтворюваність.} Артефакти згенеровано скриптом \texttt{scripts/synthetic-benchmark.ts} (детермінований, без мережі) у каталог \texttt{paper/benchmark/faculty50/}:
\begin{lstlisting}
npm run benchmark -- --researchers 50 --departments 4 --outdir paper/benchmark/faculty50
\end{lstlisting}
\textbf{Обмеження.} Коефіцієнти $a,b,c$ та параметри генерації ($N$ праць на НПП, частка дублікатів, спотворення) є явними припущеннями; їх калібрування на реальних ORCID/OpenAlex-профілях і вимірювання часу за участю користувачів --- предмет подальшої роботи.
